\chapter{Evaluating the invariant in a finite polynomial algebra}\label{ch:source}

We now turn the tangle operations of Chapter~\ref{ch:introduction} into finite matrix calculations. A strand carries a finite-dimensional vector space, a crossing acts on the two corresponding factors, and joining endpoints means contracting the appropriate indices. Closing an auxiliary circle takes a trace. Leaving one marked main strand open until the final contraction gives the reduced normalization.

The purpose of the polynomial algebra introduced below is to make these contractions compositional. Every two-strand operator that occurs here is a polynomial in one matrix, the pair Casimir $\Omega$. Stacking becomes multiplication of polynomials, while the quarter-turn has a particularly simple action in the Newton basis. The last section then expresses an arbitrary filling power in the basis of colored auxiliary-circle operators. All of these statements are finite identities at each level; asymptotics begin only in Chapter~\ref{ch:entropy}.

\section{Crossings and the one-cut normalization}
Fix $N\ge2$, and set
\begin{equation}\label{eq:root-notation}
 d=N-1,\qquad q=e^{\pi\ii/N},\qquad Q=q^2,\qquad
 x=q^d,\qquad \kappa=q^{d^2/2},\qquad \rho=q-q^{-1}.
\end{equation}
There is no restriction to prime or odd $N$. The main strands carry the $N$-dimensional vector space $V_N$; this is what ``complete color'' means at this stage. The operators $H,K,E,F$ encode its quantum-group action. They are concrete matrices, which we specify so that the crossing convention can be checked. One convenient choice of basis is $v_0,\ldots,v_d$ and
\[
 H v_j=(d-2j)v_j,\quad K=q^H,\quad
 E v_j=[j]_q^2v_{j-1},\quad Fv_j=v_{j+1},
 \qquad [j]_q=\frac{q^j-q^{-j}}{q-q^{-1}}.
\]
Missing endpoint vectors are zero. For $k\ge0$, the quantum factorial is
$[k]_q!=\prod_{r=1}^k[r]_q$, with $[0]_q!=1$.
The finite universal crossing, including its Cartan factor, is
\begin{equation}\label{eq:universal-crossing}
 \mathscr R=q^{H\otimes H/2}
 \sum_{k=0}^{d}\frac{q^{k(k-1)/2}\rho^k}{[k]_q!}
 E^k\otimes F^k.
\end{equation}
Every denominator is nonzero. The physical braided crossing is obtained with the fixed exchange and transpose convention; the inverse is the actual matrix inverse. The framing and enhancement are part of the definition, not an unspecified factor appended to an asymptotic formula.

For calculations in the polynomial algebra it is useful to change basis on this same space. We call the new coordinates balanced; they do not define a different invariant. In this basis the operators are
\begin{equation}\label{eq:balanced-module}
 \mathcal E|a\rangle=(Q^{-a}-1)|a-1\rangle,\quad
 \mathcal F|a\rangle=(1-Q^{a+1})|a+1\rangle,\quad
 K|a\rangle=xQ^{-a}|a\rangle.
\end{equation}
The ordinary cup and cap, and the one marked cap, have entries
\begin{equation}\label{eq:cups}
 p_{ab}=Nx^{-b}\delta_{a+b,d},\qquad
 \ell_{ab}=\frac{x^a}{N}\delta_{a+b,d},\qquad
 m_{ab}=\frac{x^b}{N^2}.
\end{equation}
The cup creates two endpoints and the cap joins them. Their two snake identities say that a bend followed by the opposite bend evaluates to a straight strand. The different marked cap implements the reduced invariant and satisfies $mp=1$. The diagonal operator $g=K^{1-N}=(-1)^dK$ is the \emph{pivot}, the correction inserted when an ordinary strand is closed by a trace. For a two-strand operator $A$ commuting with the pair quantum-group action, the original cylindrical evaluation is
\begin{equation}\label{eq:phi-physical}
 \phi(A)=\frac1N\Tr_{V_N\otimes V_N}\bigl[(I\otimes g)A\bigr].
\end{equation}
The factor $1/N$ in this expression is the single cut average. It is not repeated at the auxiliary circles.

\begin{proposition}[Normalization]\label{prop:normalization}
Let $A_{\rm raw}$ be the scalar obtained from the original crossings, turns, ordinary auxiliary traces and the single mark. In the enhanced reduced Kashaev convention,
\[
 |\langle L\rangle_N|=|A_{\rm raw}(L)|.
\]
In the full modified one-cut convention the modulus is $N^{-1}|A_{\rm raw}(L)|$. The changes of recorded component framing contribute only unit phases at the complete color.
\end{proposition}
\begin{proof}
The exact Fourier comparison is recorded in Appendix~\ref{app:normalization}. It is a simultaneous similarity of every adjacent braid generator and its inverse, together with an exact equality of enhancements. The factorial changes of basis cancel on every internal or traced wire. On the one remaining head the partial trace is scalar: it commutes with $E,F,K$, whose distinct weights and nonzero adjacent coefficients force scalarity. Hence the open-factor similarity cancels as well.

For clarity, the complete modified dimension is evaluated as a whole meridian germ. With the original meridian $M=-q q^\alpha=e^u$,
\[
 d_N(\alpha)=-q^{N(1-N)/2}
 \frac{q^{\alpha+1}-q^{-\alpha-1}}{q^{N\alpha}-q^{-N\alpha}}
 =\ii^{N+3}\frac{\sinh u}{\sinh Nu}.
\]
Its removable value at $u=0$ is $\ii^{N+3}/N$. Exactly one component receives this factor. The twist at completeness is $q^{-d^2/2}=\kappa^{-1}$, of modulus one. These facts prove the asserted absolute normalization without estimating any large change-of-basis matrix.
\end{proof}

\section{The full pair algebra}

An operator on two strands must commute with the joint action of the quantum-group generators if it is to represent one of our invariant tangle boxes. The formulas below specify that joint action. The subsequent lemma identifies all such commuting operators; this is the meaning of the term \emph{commutant}.
The balanced pair generators and Casimir are
\begin{equation}\label{eq:pair-casimir}
 \mathcal E_{12}=\mathcal E_1K_2+\mathcal E_2,\quad
 \mathcal F_{12}=\mathcal F_1+K_1^{-1}\mathcal F_2,\quad
 K_{12}=K_1K_2,\quad
 \Omega=\mathcal F_{12}\mathcal E_{12}+qK_{12}+q^{-1}K_{12}^{-1}.
\end{equation}
Let $\Cheb_n(z)=2T_n(z/2)$, where $T_n$ is the usual Chebyshev polynomial. Thus $\Cheb_0=2$, $\Cheb_1=z$ and $\Cheb_{n+1}=z\Cheb_n-\Cheb_{n-1}$. Define
\begin{equation}\label{eq:hermite-algebra}
 P_N(z)=\Cheb_N(z)+2=\prod_{j=0}^{d}(z-c_j),\qquad
 c_j=q^{2j+1}+q^{-2j-1},\qquad
 \cH_N=\CC[z]/(P_N).
\end{equation}
The list of roots includes multiplicity. Since $c_{d-j}=c_j$, all distinct roots are double except the simple root $-2$ at odd $N$.

\begin{lemma}[Full commutant]\label{lem:commutant}
The endomorphisms of $V_N\otimes V_N$ commuting with the three pair generators are exactly $\CC[\Omega]\simeq\cH_N$.
\end{lemma}
\begin{proof}
Let $W_T$ be the span of $|a,T-a\rangle$. The $K_{12}$ eigenvalue one occurs only on $W_d$. The matrices $\mathcal E_{12}:W_{T+1}\to W_T$ for $T<d$, and $\mathcal F_{12}:W_T\to W_{T+1}$ for $T\ge d$, have triangular maximal minors with nonzero diagonals. Consequently the entire pair is generated from $W_d$ by these raising and lowering operators.

On $W_d$, the Casimir is a symmetric tridiagonal matrix with adjacent entries $-(1-Q^{a+1})^2\ne0$. Its first basis vector is cyclic: the $k$th power has last possible coordinate $k$, with coefficient equal to the nonzero product of the first $k$ adjacent entries. An operator commuting with this block is therefore a unique polynomial of degree below $N$ in that block. Subtracting the same polynomial in $\Omega$ from a pair intertwiner gives zero on $W_d$, hence zero everywhere by generation.

The elementary quantum-shift product gives $P_N(\Omega)=0$. The cyclic block has minimal polynomial of degree $N$; its continuant, or the generic highest-weight decomposition followed by regular specialization, gives precisely \eqref{eq:hermite-algebra}. This proves the claim with all first normal directions intact.
\end{proof}

At a double root $c$, reduction modulo $(z-c)^2$ leaves
\[
 f(z)=f(c)+f'(c)(z-c)\pmod{(z-c)^2}.
\]
Thus both the value and its first derivative are independent coordinates of the operator. Throughout the paper, retaining the ``full algebra'' or the ``normal data'' means retaining this derivative as well. A polynomial can vanish at every distinct root and still define a nonzero operator in $\cH_N$. Evaluating only at roots would therefore lose part of the invariant.

\section{The actual turn and the Newton basis}

We choose a polynomial basis adapted to successive Casimir roots, including their multiplicities. Its normalization is designed so that turning a physical box reverses the basis index exactly. The same basis makes the final closure a coefficient extraction. These are the two facts that will allow us to follow every internal summation index through a large construction.
Put
\begin{equation}\label{eq:newton-basis}
 s_r=\prod_{u=1}^{r}2\sin\frac{\pi u}{N},\quad s_0=1,\qquad
 B_r(z)=\prod_{j=0}^{r-1}(c_j-z),\qquad
 \Bhat_r(z)=\frac{B_r(z)}{s_r^2},\quad 0\le r\le d.
\end{equation}
The quarter-turn is the literal contraction
\[
 \Turn(A)=(I^{\otimes2}\otimes\ell)(I\otimes A\otimes I)
 (p\otimes I^{\otimes2}).
\]
Its entries in the recorded all-primal chart are
\[
 \Turn(A)_{(a,b);(c,e)}=x^{a-e}A_{(b,d-e);(d-a,c)}.
\]
The cups and caps are those of \eqref{eq:cups}. A different source-normalized rotation can carry an additional power of $x$; throughout this manuscript $\Turn$ is the unscaled genuine turn.

\begin{proposition}[Newton turn and closure]\label{prop:newton-turn}
In the full algebra,
\begin{equation}\label{eq:newton-turn}
 \Turn\Bhat_r=\Bhat_{d-r},\qquad
 \phi\left(\sum_{r=0}^{d}\alpha_r\Bhat_r\right)=\alpha_d,
 \qquad
 \phi(\Bhat_i\Bhat_j)=\delta_{i+j,d}.
\end{equation}
\end{proposition}
The finite verification is given in Appendix~\ref{app:newton}. It uses the physical turn, its nonzero adjacent charge-block entries and the finite quantum-binomial recurrence. In particular, \eqref{eq:newton-turn} is not a generic-$q$ continuation of a formula checked only on distinct root values. The turn is linear and involutive in this basis, but is not multiplicative.

The same closure can be displayed directly in Hermite coordinates:
\begin{equation}\label{eq:hermite-trace}
 \phi(f)=(-1)^d\sum_{c\ {
m double}}
 \bigl((4-c^2)f'(c)-cf(c)\bigr)
 +\mathbf1_{N\ {
m odd}}f(-2).
\end{equation}
Equivalently,
\begin{equation}\label{eq:leading-trace}
 \phi(f)=(-1)^dN^2[z^d](f\bmod P_N).
\end{equation}
For example, at a double root $c$, writing $Q_c=P_N/(z-c)^2$ gives
$Q_c(c)=N^2/(4-c^2)$ and $Q_c'(c)/Q_c(c)=c/(4-c^2)$ from the Chebyshev differential equation. The residue of $f/P_N$ is exactly the corresponding term of \eqref{eq:hermite-trace} after multiplication by $(-1)^dN^2$. Summing residues proves \eqref{eq:leading-trace}; the simple odd root has $P_N'(-2)=N^2$. The sole shortest tridiagonal path from the first to the last coordinate has product $(-1)^dN^2$, which also identifies this scalar with the original endpoint matrix element.

\section{Positive multiplication with every normal retained}

Suppose two stacked boxes contribute basis indices $i$ and $j$. Their product is a linear combination of output indices $r$. The following coefficient formula gives both the allowed triples $(i,j,r)$ and their weights. Calling this local operation a \emph{gate} is only shorthand for this one multiplication; the geometric operation remains the vertical stack.
\begin{theorem}[Multiplication law]\label{thm:positive-product}
If $i+j<N$, then
\begin{equation}\label{eq:positive-product}
 \Bhat_i\Bhat_j=\sum_{r=\max(i,j)}^{i+j}\mu_{ijr}\Bhat_r,
 \qquad
 \mu_{ijr}=\frac{s_{i+j}s_r}
 {s_i s_j s_{r-i}s_{r-j}s_{i+j-r}}>0.
\end{equation}
If $i+j\ge N$, the product is zero in the full algebra. In particular, a product gate has precisely the integer support
\begin{equation}\label{eq:gate-support}
 i+j\le d,\qquad \max(i,j)\le r\le i+j.
\end{equation}
\end{theorem}
\begin{proof}
At a generic parameter $\zeta$, put $t=-(\zeta-\zeta^{-1})^2$. The unnormalized coefficients are
\[
 B_iB_j=\sum_r
 t^{i+j-r}\frac{[i]![j]![i+j]!}
 {[r]![r-i]![r-j]![i+j-r]!}\,B_r.
\]
Multiplication by $c_i-z$ and the identity
$(c_i-z)B_r=B_{r+1}+(c_i-c_r)B_r$, with
$c_i-c_r=t[r+i+1][r-i]$, reduce induction to
\[
 [i+1][i+j+1]=[r][r-j]+[r+i+1][i+j+1-r].
\]
This follows by expanding the four powers of $\zeta$. Endpoint coefficients follow from the same zero factors. At the principal root, all factorial indices in the case $i+j<N$ are below $N$. Since $[u]!=s_u/(2\sin(\pi/N))^u$, normalization by $s_i^2s_j^2/s_r^2$ gives \eqref{eq:positive-product}.

For $i+j\ge N$, use the exact polynomial identity $B_N/B_j=B_{N-j}$, a consequence of $c_{d-u}=c_u$. Then $B_iB_j$ is divisible by $B_N=(-1)^NP_N$ with multiplicities. This proves the zero assertion without dropping the first normal directions.
\end{proof}

\section{The literal crossing circle and its colors}

A color $j$ on an auxiliary circle means that the circle carries a $j$-dimensional representation. The main strands still carry $V_N$. We must compute the resulting operator on those two main strands, including the trace around the auxiliary circle, before replacing that circle by twists. The answer is a Chebyshev polynomial in $\Omega$, with the sign fixed by the pivot.
For $1\le j\le N$, the auxiliary module is the ordinary $j$-dimensional highest-weight module. Its loop is the actual block-braiding word $[2,1,1,2]$, with the intermediate strand dimensions retained, followed by its ordinary pivot $K_j^{1-N}$. There is no auxiliary $1/j$ or $1/N$.

\begin{proposition}[Circle identity]\label{prop:circle}
The resulting two-head operator is
\begin{equation}\label{eq:circle}
 L_j=(-1)^{j-1}U_{j-1}(\Omega/2).
\end{equation}
Here $U_0=1$, $U_1(z/2)=z$, and $U_{n+1}(z/2)=zU_n(z/2)-U_{n-1}(z/2)$.
\end{proposition}
\begin{proof}
At generic $q$, evaluate the closed monodromy on a target highest vector of weight $n$. In one direction the positive raising terms annihilate that vector; in the return direction a positive lowering term cannot return to its highest coefficient. The two Cartan factors and the auxiliary pivot give
\[
 \sum_{a=0}^{j-1}q^{(n+1)(j-1-2a)}
 =U_{j-1}\bigl((q^{n+1}+q^{-n-1})/2\bigr).
\]
The generic pair highest vectors have distinct Casimirs and descendants of total dimension $N^2$, so this proves the full matrix identity. All finite crossing denominators have indices below $N$, and the identity specializes regularly. Finally $K_j^{1-N}=(-1)^{j-1}K_j$ at the root, proving \eqref{eq:circle}, including its nilpotent coefficients.
\end{proof}

In particular $L_N=(-1)^dP_N'(\Omega)/N$. The trace identity
\begin{equation}\label{eq:circle-trace}
 \phi(L_Nf)=N\sum_{j=0}^{d}f(c_j)
\end{equation}
counts the repeated roots. It follows first by Lagrange interpolation for a polynomial with distinct roots and then by polynomial continuation in its coefficients. At even levels $L_N^2=0$; at odd levels $L_N^2=N^2E_{-2}$, with the simple-node projector $E_{-2}$. Thus complete color alone cannot be inverted to reconstruct a filling.

Let $M=\lfloor d/2\rfloor$ and $h_j=\Turn L_j$. The complete leaf has the exact positive expansion
\begin{equation}\label{eq:complete-leaf}
 h_N=\sum_{v=0}^{M}a_{N,v}\Bhat_v,
 \qquad a_{N,v}=\frac{N^2s_{2v}}{s_v^4}>0.
\end{equation}
For $0\le h\le M$, with $h_0=0$, the full color decomposition is
\begin{equation}\label{eq:core-alias}
 h_{N-h}=\sum_{v=h}^{M}
 \frac{N^2s_{2v}}{s_v^2s_{v-h}s_{v+h}}\Bhat_v+(-1)^Nh_h.
\end{equation}
For $1\le h\le\lfloor N/2\rfloor$,
\begin{equation}\label{eq:low-alias}
 h_h=\sum_{r=0}^{h-1}(-1)^{h-1+r}
 \frac{s_r^2s_{h+r}}{s_{2r+1}s_{h-r-1}}\Bhat_{d-r}.
\end{equation}
These identities separate two different ranges of Newton indices. The first term in \eqref{eq:core-alias}, called the \emph{core}, has positive coefficients before a half twist and occupies the lower indices. The second term, called the \emph{alias}, is an exact signed contribution from the complementary color; formula~\eqref{eq:low-alias} puts it in the higher Newton indices. An alias is therefore part of the finite invariant, not an error introduced by approximation. Appendix~\ref{app:colors} derives both formulas from the finite Chebyshev generating function. Chapter~\ref{ch:alias} will prove when the contribution of aliases is exponentially smaller after the complete tangle is assembled.

\section{Half twists and exact filling}

The preceding formulas describe a circle with no retained crossing. We next insert one positive or negative crossing in its original disk framing. This changes both the coefficients and their phases. We then expand a full integer twist power in the circle basis. Keeping these two steps separate ensures that the physical power is exactly $R^{\eps-2p}$, with the retained half twist counted once.
A retained half twist acts on a turned leaf by
$\Turn_\eps(f)=\Turn(R^\eps\Turn f)$, for $\eps=\pm1$. Its exact transport is
\begin{equation}\label{eq:half-twist}
 \Turn_\eps\Bhat_v=\kappa^\eps(-1)^v
 \sum_{r=0}^{v}q^{\eps v(r+1)}
 \frac{s_v}{s_rs_{v-r}}\Bhat_r.
\end{equation}
The triangular formula retains the full Cartan phase and all normals. Its derivation uses the finite continuant expansion and root factorial reflection; the two signs are coefficient conjugates in the real Newton basis. A half-twisted alias may occupy every Newton coordinate, so the pure alias support cannot be reused after \eqref{eq:half-twist}.

The $N$ circle polynomials $L_1,\ldots,L_N$ form a basis of $\cH_N$. The full crossing has a Hermite polynomial $R=\kappa r_N(\Omega)$ with
\[
 r_N(c_k)=(-1)^kq^{k(k+1)},\qquad
 \frac{r_N'(c_k)}{r_N(c_k)}
 =-\frac{d-2k}{2(q^{2k+1}-q^{-2k-1})}
\]
at each double node, and the stated value at the odd simple node. Fourier--Hermite interpolation therefore gives, for every integer $p$,
\begin{equation}\label{eq:exact-filling}
 R^{-2p}=\kappa^{-2p}\sum_{j=1}^{N}\beta_{j-1}^{-p}L_j,
 \qquad \sum_{j=1}^{N}|\beta_{j-1}^{-p}|\le2N(1+|p|).
\end{equation}
The notation $\beta_{j-1}^{-p}$ denotes the full interpolation coefficient, not a power of a scalar eigenvalue. Its value and derivative data are specified in Appendix~\ref{app:fourier}.

Insert \eqref{eq:exact-filling} at every original auxiliary disk, tensoring with the spectators before taking the unchanged closure. If $\mathcal A_N(\mathbf h)$ is the actual mixed source with colors $N-h_\ell$, then
\begin{equation}\label{eq:whole-filling-source}
 \langle K(\mathbf p)\rangle_N
 =\chi_N\kappa^{-2\sum p_\ell}
 \sum_{j_1,\ldots,j_H=1}^{N}
 \left(\prod_\ell\beta_{j_\ell-1}^{-p_\ell}\right)
 \mathcal A_N(N-j_1,\ldots,N-j_H).
\end{equation}
Here $|\chi_N|=1$ in the reduced convention. This is an all-input finite identity. The remainder of the proof analyzes this complete sum, rather than replacing it by the complete-color term.
